% GENERATED FILE. Edit canonical lessons, then run npm run generate:books.
% canonical-source-hash: fnv1a64:f37ed83930a0cc48
% canonical-lessons: PT-C317-aposentar, PT-C317-pensionista, PT-C317-emigrar, PT-C317-emigracao, PT-C317-emigrante

\chapter{Uma nova etapa}
\label{ch:pt-uma-nova-etapa}
\hlchaptermodality{317}

\begin{chapteropening}
\textbf{By the end of this chapter:} \emph{I can talk about retiring, pensioners, and leaving one’s country to live abroad.}

Complete the last lesson of chapter 317: I can talk about retiring, pensioners, and leaving one’s country to live abroad.
\end{chapteropening}

\section[aposentar]{aposentar — to retire (usually aposentar-se)}

\label{lesson:PT-C317-aposentar}



\begin{quote}
Before the new one: say the Portuguese for the will, then the Portuguese for the pension.
\end{quote}

\subsection*{You'll want to know: aposentar}
\textbf{aposentar} — \textquotedblleft{}to retire (usually aposentar-se)\textquotedblright{}.

\begin{quote}
A professora aposentou-se aos sessenta e seis anos.
\end{quote}

\begin{quote}
\emph{The teacher retired at sixty-six.}
\end{quote}

Mostly \emph{aposentar-se}, the usual word in Brazil; in Portugal \emph{reformar-se} is more common, though civil servants such as teachers \emph{aposentam-se}.

\begin{cousinweb}
\emph{aposentar} comes down from Old Portuguese \emph{apousentar} \textquotedblleft{}to lodge, to give rest\textquotedblright{} (EtymDB), built on \emph{pousar} \textquotedblleft{}to rest\textquotedblright{}, from Latin \emph{pausāre} \textquotedblleft{}to stop\textquotedblright{} (Lewis \& Short).
\end{cousinweb}

\begin{grammarlens}[title={What you've built}]
One more word for telling, explaining and describing at B1.
\end{grammarlens}

\subsection*{Your turn}
\begin{itemize}
\raggedright
  \item \emph{Say it:} \emph{aposentar}
  \item \emph{Say it:} \emph{aposentar}, once more
  \item \emph{Say it:} a sentence of your own with \emph{aposentar}
  \item \emph{Recall:} say the Portuguese for the will, then the Portuguese for the pension, then say \emph{aposentar} again
\end{itemize}

\begin{tcolorbox}[breakable,colback=teal!4,colframe=teal!35!black,title={Before you move on}]
What does aposentar mean? (\textquotedblleft{}To retire (usually aposentar-se)\textquotedblright{}.) Say it once more.
\end{tcolorbox}
\section[o pensionista]{o pensionista — the pensioner}

\label{lesson:PT-C317-pensionista}



\begin{quote}
Before the new one: say the Portuguese for the pension, then the Portuguese for to retire.
\end{quote}

\subsection*{You'll want to know: o pensionista}
\textbf{o pensionista} — \textquotedblleft{}the pensioner\textquotedblright{}.

\begin{quote}
O meu vizinho é pensionista e passa o dia no jardim.
\end{quote}

\begin{quote}
\emph{My neighbour is a pensioner and spends the day in the garden.}
\end{quote}

Built on \emph{a pensão}. The same form for a woman: \emph{a pensionista}.

\begin{cousinweb}
\emph{pensionista} is built on \emph{pensão} \textquotedblleft{}a pension\textquotedblright{} with \emph{-ista} for the person; \emph{pensão} comes from Latin \emph{pēnsiō} \textquotedblleft{}a payment\textquotedblright{} (EtymDB).
\end{cousinweb}

\begin{grammarlens}[title={What you've built}]
One more word for telling, explaining and describing at B1.
\end{grammarlens}

\subsection*{Your turn}
\begin{itemize}
\raggedright
  \item \emph{Say it:} \emph{o pensionista}
  \item \emph{Say it:} \emph{o pensionista}, once more
  \item \emph{Say it:} a sentence of your own with \emph{o pensionista}
  \item \emph{Recall:} say the Portuguese for the pension, then the Portuguese for to retire, then say \emph{o pensionista} again
\end{itemize}

\begin{tcolorbox}[breakable,colback=teal!4,colframe=teal!35!black,title={Before you move on}]
What does o pensionista mean? (\textquotedblleft{}The pensioner\textquotedblright{}.) Say it once more.
\end{tcolorbox}
\section[emigrar]{emigrar — to emigrate}

\label{lesson:PT-C317-emigrar}



\begin{quote}
Before the new one: say the Portuguese for to retire, then the Portuguese for the pensioner.
\end{quote}

\subsection*{You'll want to know: emigrar}
\textbf{emigrar} — \textquotedblleft{}to emigrate\textquotedblright{}.

\begin{quote}
Nos anos sessenta, o meu avô emigrou para o estrangeiro.
\end{quote}

\begin{quote}
\emph{In the sixties my grandfather emigrated abroad.}
\end{quote}

\emph{Emigrar} is leaving; \emph{imigrar} is arriving.

\begin{cousinweb}
\emph{emigrar} is a learned borrowing from Latin \emph{ēmigrāre} \textquotedblleft{}to move out\textquotedblright{}, built on \emph{migrāre} \textquotedblleft{}to move from place to place\textquotedblright{} (Lewis \& Short).
\end{cousinweb}

\begin{grammarlens}[title={What you've built}]
One more word for telling, explaining and describing at B1.
\end{grammarlens}

\subsection*{Your turn}
\begin{itemize}
\raggedright
  \item \emph{Say it:} \emph{emigrar}
  \item \emph{Say it:} \emph{emigrar}, once more
  \item \emph{Say it:} a sentence of your own with \emph{emigrar}
  \item \emph{Recall:} say the Portuguese for to retire, then the Portuguese for the pensioner, then say \emph{emigrar} again
\end{itemize}

\begin{tcolorbox}[breakable,colback=teal!4,colframe=teal!35!black,title={Before you move on}]
What does emigrar mean? (\textquotedblleft{}To emigrate\textquotedblright{}.) Say it once more.
\end{tcolorbox}
\section[a emigração]{a emigração — emigration}

\label{lesson:PT-C317-emigracao}



\begin{quote}
Before the new one: say the Portuguese for the pensioner, then the Portuguese for to emigrate.
\end{quote}

\subsection*{You'll want to know: a emigração}
\textbf{a emigração} — \textquotedblleft{}emigration\textquotedblright{}.

\begin{quote}
Nos anos sessenta, a emigração esvaziou muitas aldeias.
\end{quote}

\begin{quote}
\emph{In the sixties emigration emptied many villages.}
\end{quote}

Leaving the country; the opposite movement is immigration.

\begin{cousinweb}
\emph{emigração} is a learned word from Late Latin \emph{ēmigrātiō} \textquotedblleft{}a moving out\textquotedblright{}, built on \emph{migrāre} \textquotedblleft{}to move from place to place\textquotedblright{} (Lewis \& Short).
\end{cousinweb}

\begin{grammarlens}[title={What you've built}]
One more word for telling, explaining and describing at B1.
\end{grammarlens}

\subsection*{Your turn}
\begin{itemize}
\raggedright
  \item \emph{Say it:} \emph{a emigração}
  \item \emph{Say it:} \emph{a emigração}, once more
  \item \emph{Say it:} a sentence of your own with \emph{a emigração}
  \item \emph{Recall:} say the Portuguese for the pensioner, then the Portuguese for to emigrate, then say \emph{a emigração} again
\end{itemize}

\begin{tcolorbox}[breakable,colback=teal!4,colframe=teal!35!black,title={Before you move on}]
What does a emigração mean? (\textquotedblleft{}Emigration\textquotedblright{}.) Say it once more.
\end{tcolorbox}
\section[o emigrante]{o emigrante — the emigrant}

\label{lesson:PT-C317-emigrante}



\begin{quote}
Before the new one: say the Portuguese for to emigrate, then the Portuguese for emigration.
\end{quote}

\subsection*{You'll want to know: o emigrante}
\textbf{o emigrante} — \textquotedblleft{}the emigrant\textquotedblright{}.

\begin{quote}
O emigrante volta à aldeia todos os verões.
\end{quote}

\begin{quote}
\emph{The emigrant comes back to the village every summer.}
\end{quote}

The same form for a woman: \emph{a emigrante}.

\begin{cousinweb}
\emph{emigrante} is a learned borrowing from Latin \emph{ēmigrāns} \textquotedblleft{}moving out\textquotedblright{} (EtymDB), from \emph{ēmigrāre}, built on \emph{migrāre} \textquotedblleft{}to move\textquotedblright{}.
\end{cousinweb}

\begin{grammarlens}[title={What you've built}]
One more word for telling, explaining and describing at B1.
\end{grammarlens}

\subsection*{Your turn}
\begin{itemize}
\raggedright
  \item \emph{Say it:} \emph{o emigrante}
  \item \emph{Say it:} \emph{o emigrante}, once more
  \item \emph{Say it:} a sentence of your own with \emph{o emigrante}
  \item \emph{Recall:} say the Portuguese for to emigrate, then the Portuguese for emigration, then say \emph{o emigrante} again
\end{itemize}

\begin{tcolorbox}[breakable,colback=teal!4,colframe=teal!35!black,title={Before you move on}]
What does o emigrante mean? (\textquotedblleft{}The emigrant\textquotedblright{}.) Say it once more.
\end{tcolorbox}
